% Reflection essay — 500-700 words.
% NOTE: items marked [PERSONAL: ...] MUST be rewritten from your real
% experience (or deleted) before submission. Do not submit with markers.
When I started this project, the goal was not to build the biggest library
system in the world, but to walk through a complete software engineering
process on a system small enough to understand end-to-end. The Smart Library
\& Fine Management System uses the MERN stack --- React with Tailwind CSS on
the front, Express and Mongoose with MongoDB behind a REST API --- because it
let me keep one language (JavaScript) across the whole project and focus on
engineering practice rather than fighting tooling.

For the process model I chose the Incremental Model. The five modules ---
members, books, borrow/return, fines and audit --- are naturally separable,
so I could finish and test each slice before starting the next. This suited
a single-developer sessional project: if something went wrong in the fine
policy, member registration was already stable and demonstrable.
[PERSONAL: describe one real problem you hit in a specific increment and how
incremental delivery helped.]

Three design patterns did real work in the code. The \textbf{Singleton}
\texttt{DatabaseManager} gives the process one shared connection manager, so
the connection lifecycle has a single owner; I rejected Prototype because
cloning a database manager would be meaningless --- there is nothing useful
to copy. The \textbf{Facade} \texttt{BorrowFacade} hides the multi-step
issue/return workflow (member lookup, borrow limit, stock update, due date,
fine calculation, audit logging) behind two readable methods; I rejected
Adapter because no incompatible third-party interface needed translating.
The \textbf{Strategy} pattern implements the fine policies --- 10~BDT/day
for students, zero for faculty, and a heavy-overdue decorator charging
20~BDT/day after the first seven days. I rejected State because the fine
rules are interchangeable algorithms selected by role, not an object whose
behaviour changes with internal state.

Test-Driven Development was applied most strictly to fine calculation. I
recorded a real Red--Green cycle: the tests failed against the unimplemented
strategy (``five overdue days $\rightarrow$ 50~BDT'' expected 50, received
0), passed once the ten-taka-per-day rule was implemented, and the boundary
cases (day~7~$=$~70, day~8~$=$~90) were locked in with further tests.
Mocking mattered in the facade unit tests: repositories, the fine service
and the audit service are all Jest mocks, so the tests verify coordination
logic without a database and run in milliseconds. Integration tests with
Supertest cover the real HTTP path against an isolated in-memory MongoDB,
which catches wiring mistakes mocks cannot see. The Selenium test drives the
actual React UI in Chrome --- add book, add member, issue, return --- and
passed on my machine.
[PERSONAL: describe one real bug that Selenium or the integration tests
found that unit tests missed.]

Continuous integration through GitHub Actions re-runs tests, coverage and
the frontend build on every push, which turns ``it works on my machine''
into a shared, visible fact. As the self-directed tool I learned
\textbf{Swagger/OpenAPI}: one YAML file now generates interactive
documentation at \texttt{/api-docs}, and I can test any endpoint from the
browser. I chose it over a README because a specification is structured,
versionable and much harder to let drift out of date.
[PERSONAL: mention which Swagger resources you actually used and one real
difficulty, e.g.\ YAML indentation.]

The hardest part was [PERSONAL: replace with your real challenge --- e.g.\
understanding atomic MongoDB updates, React select state in Selenium, or
wiring dependency injection]. I resolved it by [PERSONAL: your real solution
steps]. The learning resources I relied on were mostly official
documentation (Mongoose, Jest, Express, Tailwind), which I evaluated by
whether they matched the project's library versions and whether I could
apply them within one sitting.

The system has clear limitations: no authentication, no payment flow, regex
search instead of Atlas Search, and no multi-document transactions --- the
issue workflow instead uses an atomic conditional stock update, which I
documented honestly rather than overclaiming. Next I want to learn proper
authentication (JWT with role checks) and containerised deployment, because
the realistic next step for this application is running it safely for real
users rather than adding more features.

[FINAL STEP: after your viva rehearsals and real test runs, re-read the
essay aloud, keep it between 500 and 700 words, and remove every bracketed
marker.]
